\section{Research questions and design}
\label{sec:design}

\subsection{The question the design is built around}

The exposure that preceded SVB's failure was not hidden. Some of it was published,
in the filing reproduced in Table~\ref{tab:disclosed}, and the supervisory record
indicates that the direction of the problem was recognised before it became acute
\citep{FRB2023}. What was not published was any economic-value measure at the
point it would have mattered most, since the 2022 filing dropped the profile that
the 2021 filing had carried.

That gap suggests a question that can be answered rather than argued about. An
analyst at the end of 2021, holding nothing but what SVB had already filed and
freely available market data, would have had a disclosed economic-value profile
to work with. Could such an analyst have built a measure of the exposure that
survived contact with the following year? And would the same measure, applied
without modification to a bank that did not fail, have distinguished between
them?

This gives three questions.

\begin{enumerate}
\item Can a valuation model calibrated only on information available at
31~December 2021 reproduce the economic-value profile that SVB published for that
date, at shock points it was not fitted to?
\item Applied identically to both banks and with its spreads held fixed, does the
same model price their securities portfolios a year later, against what the banks
subsequently disclosed?
\item What assumption about depositor behaviour is required to make SVB's own
published sensitivity internally consistent, and is that assumption supportable
for the depositor base SVB actually had?
\end{enumerate}

The third is where the paper ends up, and it is worth saying why at the outset.
An economic-value-of-equity figure is not a measurement in the way a securities
fair value is. It is the net of an asset valuation, which the filings constrain
tightly, and a deposit valuation, which they barely constrain at all. The
asset side can be rebuilt and checked. The liability side depends on how long
deposits are assumed to stay, and that assumption is not disclosed, not
observable before the fact, and capable of moving the answer by more than the
entire asset-side exposure.

\subsection{Design}

The model is calibrated at one date and tested at another, with nothing refitted
in between.

Calibration uses 31~December 2021 only. Securities are valued off a fitted zero
curve with a short-rate model behind it, so that positions carrying prepayment
options can be priced at all, and one option-adjusted spread per security-type and
maturity bucket is solved so the model reproduces each bank's disclosed fair values
at that date. Non-maturity deposits are valued as a below-market funding source
that decays, following \citet{Drechsler2021}, in which deposits hedge a
long-duration asset book precisely to the extent that they are sticky and
insensitive to rates. Shocks are the standardised scenarios of the Basel Committee's
2016 standard, in the USD calibration in force at both valuation dates, with the
supervisory outlier test set at 15\% of Tier~1 capital \citep{BCBS2016}.
Section~\ref{sec:reconstruction} gives each of these in full; what matters for the
design is what happens next.

The spreads are then frozen. The test is to re-price both portfolios at
31~December 2022 under the curve that actually obtained, ageing each position and
changing nothing else, and to compare the result against the fair values the banks
subsequently disclosed. No parameter is permitted to see the 2022 data. This is
the subject of Section~\ref{sec:comparison}.

Of the parameters the construction needs, almost all are either disclosed or
estimable from public series. Two are neither. The deposit beta --- the pass-through
from market rates to deposit rates --- is disclosed by SVB at 0.60, so it is used
rather than estimated. The half-life of the deposit balance is disclosed by nobody,
and it is the pivot of the whole exercise.

\subsection{What is fitted and what is not}

One parameter is solved against the disclosed profile. Holding beta at 0.60, the
uninsured-deposit half-life is chosen so that the model reproduces the $-13.9\%$
that SVB reported at $+100$ basis points. That solve returns a half-life of
3.59~years for uninsured deposits and 12.56~years for insured.

Everything else in the profile is then out of sample in the ordinary sense: the
model was not shown it, and nothing was adjusted to improve the fit to it. The
$+100$ basis point point is a calibration target and is reported as such
throughout. It is not evidence that the model works, and no agreement at that
point is claimed. The $+200$ basis point point, both downward shocks, and the
base level of economic value are the places where the model can be wrong, and
Section~\ref{sec:results} reports what it does there.

\subsection{What the design can and cannot establish}

Three boundaries are worth fixing before any results appear, because the
temptation to read past each of them is real.

First, the results are deterministic. The model is solved, not estimated from a
sample, and the uncertainty reported in this paper is sensitivity to assumptions
rather than sampling variation. Section~\ref{sec:sensitivity} reports ranges
across explicit parameter variants. No standard errors, confidence intervals or
significance tests are offered for any economic quantity, and the ranges should
not be read as though they were. The one exception is the Monte Carlo error of the
pricing engine itself, which is reported in Appendix~\ref{app:pricing} and
concerns the numerical precision of a simulation whose answer is known
analytically.

Second, the paper makes no causal claim about SVB's failure. The design can
establish whether the exposure was measurable from what had been published. It
cannot establish that disclosing an economic-value figure in the 2022 filing
would have changed the outcome, and no such counterfactual is attempted. Where
the paper moves from measurement to reading, Section~\ref{sec:interpretation}
says so.

Third, the economic-value figures computed for Schwab have no disclosed
comparator, because Schwab published none. They are model output under a stated
deposit assumption and are labelled that way wherever they appear. They are
included because applying one framework to both banks is what makes the
comparison informative, not because they are measurements of Schwab's exposure.
